\section{Related work}

\paragraph{DOA with small arrays.} GCC-PHAT \cite{knapp1976gcc} and
steered-response power with PHAT weighting (SRP-PHAT)
\cite{dibiase2001srp} are the standard robust estimators under
reverberation; ODAS \cite{grondin2019odas} runs SRP-PHAT with tracking on
the same ReSpeaker and similar arrays for robots. Speech-presence or
SNR-dependent weighting of time--frequency bins is a known improvement for
GCC-type methods in noise; we evaluate it, and noise cross-spectrum
subtraction, as two uses of one silence model.

\paragraph{Distance estimation.} Source distance from a compact array
relies on the direct-to-reverberant ratio (DRR), level, or learned
features \cite{vesa2009binaural}. The coherent-to-diffuse ratio (CDR) is a
coherence-based DRR proxy with closed-form estimators
\cite{schwarz2015cdr}; it depends on the room's critical distance, which
scales with $\sqrt{V/\mathrm{RT}_{60}}$ \cite{kuttruff2016room}.

\paragraph{Spatial filtering before enhancement.} MVDR beamforming
\cite{capon1969mvdr} followed by single-channel enhancement is the classic
multichannel speech-enhancement chain; small-aperture arrays need
superdirective designs and are sensitive to steering errors. Neural
single-channel denoisers such as DeepFilterNet3
\cite{schroeter2023deepfilternet3} remove stationary and diffuse noise but
cannot separate two talkers by direction.

\paragraph{Simulation.} We use the image-source method
\cite{allen1979image} via pyroomacoustics \cite{scheibler2018pra}, and a
diffuse noise field with the exact spherical-isotropic coherence
\cite{habets2008diffuse}.
